% LuaLaTeX / luatexja 用の Beamer 日本語スライド。
\documentclass{beamer}
\usepackage{luatexja}

\usetheme{default}
\title{人工知能による構文解析とプレゼンテーション生成}
\subtitle{ロスレス具体構文木(CST)と非破壊的編集の実現}
\author{研究開発チーム}
\date{2026年10月}

\begin{document}

\begin{frame}
  \titlepage
\end{frame}

\section{研究の背景と課題}

\begin{frame}{研究背景と既存手法の限界}
  \framesubtitle{全角記号「」や約物、※印を含む文書構造の検証}
  \begin{itemize}
    \item 従来のWYSIWYGツールでは、Gitによる版管理が困難。
    \item LaTeXソースを直接編集する際、意図せぬコード破壊(コメント消失等)が発生しやすい。
    \item 空白・改行・マルチバイト文字の1バイト単位での厳密な保持。
  \end{itemize}

  \begin{block}{提案手法の核となる着眼点}
    中粒度アイランドパースを採用し、構造化対象外のマクロを安全に隔離する。
  \end{block}
\end{frame}

\end{document}
